\documentclass[a4paper,11pt]{article}
\usepackage[utf8]{inputenc}
\usepackage[T1]{fontenc}
\usepackage[french]{babel}
\usepackage[left=1cm,right=1cm,top=.5cm,bottom=0cm]{geometry}
\usepackage{enumitem}
\usepackage{hyperref}
\usepackage{xcolor}
\usepackage{titlesec}
\usepackage{fontawesome5}
\usepackage{lmodern}
\usepackage[normalem]{ulem}

\definecolor{primary}{RGB}{0, 51, 102}
\definecolor{secondary}{RGB}{80, 80, 80}

\hypersetup{colorlinks=true, linkcolor=primary, filecolor=primary, urlcolor=primary}
\titleformat{\section}{\large\bfseries\color{primary}\uppercase}{}{0em}{}[\titlerule]
\titlespacing{\section}{0pt}{8pt}{5pt}

\begin{document}
\begin{center}
    {\Large \textbf{Lo\"ic Aron MBASSI EWOLO}} \\ \vspace{4pt}
    \textbf{\large Alternant Ing\'enieur IA Agentique -- Syst\`emes Intelligents} \\
    \textit{\small Stage INRIA Saclay (Avril--Ao\^ut 2026) -- Disponible en alternance d\`es septembre 2026 (12 mois)} \\ \vspace{4pt}
    \small
    \faMapMarker\ Cergy, France (Mobile IDF) \quad \faPhone\ (+33) 7 59 52 33 98 \quad
    {\color{primary}\faEnvelope\ \href{mailto:loicaron.mbassiewolo@ensea.fr}{loicaron.mbassiewolo@ensea.fr}} \\ \vspace{2pt}
    {\color{primary}\faLinkedin\ \href{https://www.linkedin.com/in/loic-mbassi/}{linkedin.com/in/loic-mbassi}} \quad
    {\color{primary}\faGithub\ \href{https://github.com/Nameless0l}{github.com/Nameless0l}} \quad
    {\color{primary}\faGlobe\ \href{https://mbassiloic.tech}{mbassiloic.tech}}
\end{center}
\vspace{-6pt}
\noindent\small{Architecture d'un syst\`eme multi-agents (Google ADK, RAG, LangChain) avec orchestration et résolution de conflits entre agents spécialisés, et déploiement d'un pipeline NLP (BERT, scoring RGPD) à l'INRIA Saclay. Double diplôme ENSEA/Polytechnique Yaoundé en IA et mathématiques appliquées. Expérience en RAG intégré dans Snappy (Gemini Flash/Pro) et recherche sur la généralisation des LLM (MILA Montréal). BPCE SI, moteur technologique du groupe BPCE, est un cadre idéal pour déployer des architectures agentiques intelligentes au c\oe{}ur de services bancaires numériques.}

\section{Formation}
\begin{itemize}[leftmargin=*, label=\tiny$\bullet$, nosep]
    \item \textbf{ENSEA -- \'Ecole Nationale Sup\'erieure de l'\'Electronique et de ses Applications} \hfill \textit{2025 -- 2027} \\
    \small{Double dipl\^ome d'Ing\'enieur -- Syst\`emes \'Electroniques, Signal, IA \& Cybersécurité.}
    \item \textbf{\'Ecole Polytechnique de Yaound\'e (1\textsuperscript{\`ere} \'ecole d'ing\'enieur CEMAC)} \hfill \textit{2021 -- 2025} \\
    \small{Dipl\^ome d'Ing\'enieur de Conception (Master 2) : \textbf{Math\'ematiques Appliqu\'ees}, G\'enie Logiciel, Algorithmique.}
\end{itemize}

\section{Comp\'etences Techniques}
\begin{itemize}[leftmargin=*, nosep]
    \item \small{\textbf{IA Agentique :} Google ADK, LangChain, RAG, orchestration multi-agents, raisonnement collaboratif, résolution conflits, planning multi-étapes, agents spécialisés.}
    \item \small{\textbf{NLP / LLM :} BERT, Transformers, fine-tuning, HuggingFace, extraction d'entités, classification, scoring automatisé, interprétabilité.}
    \item \small{\textbf{Backend / APIs :} Spring Boot 3, Laravel, PostgreSQL, Redis, Docker, REST APIs, gestion sessions, authentification, monitoring.}
    \item \small{\textbf{MLOps :} FastAPI, pipelines CI/CD, experiments tracking, monitoring modèles, gestion versions, déploiement production.}
    \item \small{\textbf{Langues :} Français (natif), Anglais (professionnel, rédaction technique EN), Camerounais.}
\end{itemize}

\section{Exp\'eriences \& Projets}
\begin{itemize}[leftmargin=*, label={-}]
    \item \textbf{Stage INRIA Saclay -- \'Equipe TRIBE (Projet PRIMO)} \hfill \textit{Avril -- Ao\^ut 2026} \\
    \textit{Palaiseau (IP Paris)} \hfill \textit{Python, NLP, BERT, Transformers, scikit-learn, pandas, fine-tuning}
    \vspace{-2pt}
    \begin{itemize}[leftmargin=0.4cm, nosep, label=\tiny$\bullet$]
        \item \small{Fine-tuning BERT/Transformers pour classification automatique de clauses de confidentialité à large échelle (500+ SDKs mobiles) ; scoring interprétable RGPD.}
        \item \small{Pipeline NLP complet : extraction politiques, tokenization, fine-tuning avec hyperparameter tuning, interprétabilité (LIME, attention weights).}
        \item \small{Documentation méthodologie et résultats ; collaboration LVMT/École des Ponts (programme PEPR MOBIDEC, financement ANR).}
        \item \small{Intégration résultats dans architecture gouvernance de conformité pour équipes non-techniques.}
    \end{itemize}
    \vspace{3pt}

    \item \textbf{Syst\`eme Multi-Agents Agricole} \hfill \textit{2024 -- 2025} \\
    \textit{Orchestration Agentique \& RAG} \hfill \textit{Google ADK, RAG, LangChain, Gemini Flash/Pro, FastAPI, Docker}
    \vspace{-2pt}
    \begin{itemize}[leftmargin=0.4cm, nosep, label=\tiny$\bullet$]
        \item \small{Architecture 4 agents spécialisés (météo, ressources, marchés, risques) orchestrés via Google ADK : raisonnement collaboratif, planification multi-étapes, résolution conflits d'objectifs.}
        \item \small{Gemini LLM pour raisonnement d'agents ; RAG avec enrichissement contextuel (APIs OpenWeather, données marchés) ; reduction hallucinations via retrieval augmentation.}
        \item \small{Déploiement FastAPI avec gestion erreurs, logging, monitoring ; tests robustesse ; recommandations multi-critères (rendement, coûts, risques) ; scalabilité.}
        \item \small{Documentation architecture agentique ; généralisation à autres domaines (supply chain, consulting, aide à la décision).}
    \end{itemize}
    \vspace{3pt}

    \item \textbf{Plateforme Snappy -- RAG \& Chatbots Agents} \hfill \textit{2024 -- 2025} \\
    \textit{RAG \& Orchestration LLM} \hfill \textit{Spring Boot 3, Gemini, RAG, LangChain, vector store, PostgreSQL, Redis}
    \vspace{-2pt}
    \begin{itemize}[leftmargin=0.4cm, nosep, label=\tiny$\bullet$]
        \item \small{Intégration Gemini Flash/Pro dans Snappy : chatbots intelligents orchestrant retrieval contextuel, raisonnement, génération de réponses.}
        \item \small{Vector store pour retrieval semantique (documents utilisateurs, FAQ, historique) ; fine-tuning léger pour domaine métier.}
        \item \small{Gestion contexte long (conversations multi-turn), fallback strategies, monitoring qualité ; latency optimization (batching, caching).}
        \item \small{Déploiement production avec logs structurés ; documentation système LLM pour équipes maintenance.}
    \end{itemize}
    \vspace{3pt}

    \item \textbf{Dev Backend -- Fintech CSC} \hfill \textit{Oct. 2024 -- Jan. 2025} \\
    \textit{Sécurité Transactions \& APIs} \hfill \textit{Laravel, MySQL, REST APIs, JWT, rate limiting, Docker}
    \vspace{-2pt}
    \begin{itemize}[leftmargin=0.4cm, nosep, label=\tiny$\bullet$]
        \item \small{Architecture backend transactions : authentification JWT, gestion sessions, rate limiting, prévention injections, conformité PCI-DSS.}
    \end{itemize}
    \vspace{3pt}

    \item \textbf{Assistant de Recherche -- MILA Montréal} \hfill \textit{Juin -- Sept. 2025} \\
    \textit{Théorie LLM \& Généralisation} \hfill \textit{Python, PyTorch, TensorFlow, probabilités, mathématiques}
    \vspace{-2pt}
    \begin{itemize}[leftmargin=0.4cm, nosep, label=\tiny$\bullet$]
        \item \small{Étude phénomène grokking (généralisation tardive LLM) : reproduction SOTA, analyses mathématiques, pipelines tests statistiques rigoureux.}
    \end{itemize}
    \vspace{3pt}

    \item \textbf{Laravel API Generator -- Open Source} \hfill \textit{2024 -- présent} \\
    \textit{Génération Code Backend} \hfill \textit{PHP, Laravel, Python, LLM, XML, Packagist}
    \vspace{-2pt}
    \begin{itemize}[leftmargin=0.4cm, nosep, label=\tiny$\bullet$]
        \item \small{Génération code backend Laravel depuis UML/XML via LLM ; +30 étoiles GitHub, Packagist (~20 téléchargements) ; orchestration LLM pour generation multi-étapes.}
    \end{itemize}
    \vspace{3pt}

\end{itemize}

\section{Distinctions \& Leadership}
\begin{itemize}[leftmargin=*, label=\tiny$\bullet$, nosep]
    \item \small{\textbf{1\textsuperscript{er} prix IndabaX Africa 2025} (IA Santé, hackathon panafricain) -- classification maladies, déploiement API Flask}
    \item \small{\textbf{1\textsuperscript{er} prix Mountain Hub Régional 2024} (Innovation technologique)}
    \item \small{\textbf{1\textsuperscript{er} prix CITS National Hackathon} (Smart City, ML) -- vs 75 équipes}
    \item \small{\textbf{Head of Projects Cell \& Lead AI Cell ENSEA} -- structuration 3 pôles techniques, +50\% membres actifs, collaborations Google DeepMind}
    \item \small{\textbf{Open Source :} Laravel API Generator (+30 étoiles GitHub, Packagist) -- génération code API via LLM+UML}
    \item \small{\textbf{Certifications :} Supervised ML (Stanford/DeepLearning.AI/Coursera, Oct. 2024), Python for Data Science (IBM, Jan. 2024)}
\end{itemize}

\end{document}
